\emph{(i).} The equivalence of real chirality with $n \equiv 2 \pmod 4$ and with residues $0$ and $4$ among causal signatures is Theorem \claimref{synthesis.oriented-causal-towers-mirror-pair-every-window}(ii) and (iii) (with Lemma \claimref{transport.sign} for the sign of $I^{2}$): causal signatures of even dimension $n$ have residue $\pm(n-2) \pmod 8$, so $n \equiv 2 \pmod 4$ gives $r \in \{0, 4\}$ with $I^{2} = +1$ and $n \equiv 0 \pmod 4$ gives $r \in \{2, 6\}$ with $I^{2} = -1$, and for even $n$ the pseudoscalar is central in $\Cl^{0}(V)$; for odd $n$, $\Cl^{0}(V)$ is central simple. The form of the splitting on a single ideal, that reversion exchanges the two chiral projectors exactly when $n \equiv 2 \pmod 4$, is the lemma of the result filed alongside (synthesis.reversion-pairing-vanishes-residues-zero-four), and is not reproved here. What is added is the eigenvalue on the ideals of Theorem \claimref{dark-tower-matter.constituents-field} and of the lowest-causal content result (\claimref{synthesis.lowest-causal-residues-zero-four-content}). At $(1,1)$, $I = e_{t}e_{a}$, so $If = \tfrac12(I + I^{2}) = f$ and $I\psi = Iaf = aIf = \psi$ for $\psi = af$, $a \in \Cl^{0}(V)$. At $(1,5)$ and $(5,1)$, $I(1 + cI) = c(1 + cI)$ gives $If = cf$ and, $I$ being central in $\Cl^{0}(V)$, $I\psi = c\psi$ on $S$. A $\Cl^{0}(V)$-module map commutes with the central element $I$, so it preserves its eigenvalue: the ideals with $c = 1$ and $c = -1$ are not isomorphic, and the same holds at $(1,1)$ for $f$ and $1 - f$. Any conjugation of the constituent field acts on $S$ by right multiplication or by a real-linear map commuting with the action of $\Cl^{0}(V)$, so it keeps a real eigenvalue of $I$: in a really chiral tower a constituent and its conjugate have the same chirality. At residues $2$ and $6$, $I^{2} = -1$, $I$ has no real eigenvector, and chirality is the splitting into $\pm i$ for the complex structure $I$, which a conjugation reverses. The check enumerates every causal signature with $n \le 14$ and builds $S$ at $(1,1)$, $(1,5)$ and $(5,1)$.

\emph{(ii).} With $y = x^{2}$, $(x/2)/\sinh(x/2) = 1/\sum_{k} (y/4)^{k}/(2k+1)! = 1 - y/24 + 7y^{2}/5760 - 31y^{3}/967680 + 127y^{4}/154828800 + \cdots$; the product over the Chern roots, rewritten in the elementary symmetric functions $p_{j}$ of the $y_{i}$, gives the degree-$4j$ parts $-p_{1}/24$, $(7p_{1}^{2} - 4p_{2})/5760$, $(-31p_{1}^{3} + 44p_{1}p_{2} - 16p_{3})/967680$ and $(381p_{1}^{4} - 904p_{1}^{2}p_{2} + 208p_{2}^{2} + 512p_{1}p_{3} - 192p_{4})/464486400$, all nonzero (the check expands the product exactly in four variables). The anomaly theorem, part of (A), imported from the literature, cited from memory and not in the store (Alvarez-Gaum\'e and Witten, 1984), states that the local rotation-gauge and gauge anomalies of a chiral spin-$\tfrac12$ field in dimension $n$ descend from the degree-$(n+2)$ part of $\hat{A}(R)\,\mathrm{ch}(F)$, with sign set by the chirality, and that a local anomaly is removable by a local counterterm if and only if that form vanishes. Pontryagin classes have degrees $4j$, so the pure rotation-gauge part vanishes for $n \equiv 0 \pmod 4$ and is the nonzero $\hat{A}_{(n+2)/4}$ for $n \equiv 2 \pmod 4$. The pure-gauge part is $\mathrm{tr}_{R}F^{k}/k!$ with $k = (n+2)/2$, even exactly when $n \equiv 2 \pmod 4$. For a Hermitian generator $X$ of the gauge algebra with real weights $\lambda_{i}$ on $R$, $\mathrm{tr}_{R}X^{k} = \sum_{i}\lambda_{i}^{k}$; for even $k$ the conjugate representation gives the same sum and the sum is positive whenever $X$ acts nontrivially, while for odd $k$ the conjugate gives the negative. By (i) every field of a really chiral tower's single-ideal content, conjugates included, has the same chirality, so the forms of the content are sums of terms of one sign: the rotation-gauge part is $N\hat{A}_{(n+2)/4}$ for $N$ field components counted in units of one chiral field, nonzero for $N \ge 1$, and the gauge part is positive on every nontrivially acting generator. Gauge vectors are not chiral and contribute nothing; the series has no gravitino, gravity being transport and not a quantised field, and the constituents' content contains no two-form field, self-dual or not; a Green-Schwarz term, which could cancel a factorising part, needs such a field, so no field of the content can cancel these forms. The series already treats such conditions as necessary for consistency (Fact \claimref{standard-model.imported-anomaly-coefficients}, the four-dimensional case, where conjugation reverses chirality and the cubic sums can and do cancel).

\emph{(iii).} At $(1,1)$ both generators are shared, and under the sourcing rule (Lemma \claimref{dark-matter.projection-every-towers-content-dilutes-ones}) the tower's matter sits in comoving copies of its plane in this tower's spacetime and moves in the field of the one metric. A coordinate plane $y = z = 0$ of a spatially flat expanding field $a(\eta)^{2}(-d\eta^{2} + dx^{2} + dy^{2} + dz^{2})$ is the fixed set of two reflections, hence totally geodesic, with induced field $a^{2}(-d\eta^{2} + dx^{2})$ and scalar curvature $R = 2(aa'' - a'^{2})/a^{4}$ (sign convention with de Sitter positive; the store's opposite overall sign does not affect whether $R$ varies). For $a = \eta$ (radiation) $R = -2\eta^{-4}$, for $a = \eta^{2}$ (matter) $R = -4\eta^{-6}$, both with nonzero $\partial_{\eta}R$; for $a = -1/(H\eta)$ (de Sitter) $R = 2H^{2}$. The $t$-$r$ plane of the Schwarzschild field, fixed by the reflections of the sphere, has field $-f\,dt^{2} + dr^{2}/f$, $f = 1 - 2M/r$, and $R = -f'' = 4M/r^{3}$, $\partial_{r}R = -12M/r^{4} \ne 0$. The check computes these exactly. The single ideal at $(1,1)$ is one-dimensional and real with $I = +1$, so $N$ constituent fields are $N$ real chiral fields of one sign, with gravitational anomaly coefficient $c_{L} - c_{R} = \pm N/2$ (one such field has central charge $\tfrac12$). The two-dimensional anomaly in covariant form, part of (A) (imported, cited from memory: Alvarez-Gaum\'e and Witten 1984 for the anomaly, Iso, Umetsu and Wilczek 2006 for the covariant form; Robinson and Wilczek 2005 used the consistent form), is $\nabla_{\mu}T^{\mu\nu} = \kappa\,(c_{L} - c_{R})\,\epsilon^{\nu\mu}\partial_{\mu}R$ with $\kappa = \pm 1/96\pi$ by convention, and by Bardeen and Zumino the local counterterm that restores conservation leaves an antisymmetric stress proportional to $R$, so no choice of counterterm gives a stress that is both conserved and rotation-covariant. On the plane fields above $\partial R \ne 0$, so the stress of single-ideal residue-$0$ matter, massive or massless, is not conserved. The field equations of the one metric equate a divergence-free side to the source (Proposition \claimref{dark-matter.sourcing-rule-forced-not-chosen}), so such matter cannot be a source of them, and P1 makes every tower's matter a source. Inflow from another tower could cancel the anomaly only through a coupling between towers, and the towers interact only through the one metric (Fact \claimref{dark-tower-matter.imported-towers-their-types}). The argument concerns the tower's single-ideal matter only. Content carried with equal chiral weight on the two sheets of an edge, as the horizon modes are when labelled by null direction (Theorem \claimref{synthesis.r-loc-needs-sheet-labelled-modes}(v)), has $c_{L} = c_{R}$ on the $t$-$r$ plane and no anomaly; nothing here constrains it. The argument also needs both directions of the tower to be shared, which holds at $(1,1)$ only: a higher causal member of residue $0$, $(1,9)$ or $(9,1)$, shares four directions, and against the one metric its content is four-dimensional as in (iv), with no rotation-gauge anomaly. With the content result that no number of single-ideal constituents at $(1,1)$ is massive (\claimref{synthesis.lowest-causal-residues-zero-four-content}(ii)), residue $0$ at its lowest causal member holds neither single-ideal matter nor single-ideal radiation, given (A) and the sourcing reading.

\emph{(iv).} At $(1,5)$ the shared generators are $e_{0}, \ldots, e_{3}$ and the unshared $e_{4}, e_{5}$; at $(5,1)$ they are the time $e_{5}$ with $e_{0}, e_{1}, e_{2}$, and $e_{3}, e_{4}$. In both orderings $I = I_{4}J$ with $I_{4}$ the ordered product of the shared generators and $J$ that of the unshared ones, $I_{4}^{2} = J^{2} = -1$, and $I_{4}$, $J$ commute. On $S$, $I_{4}J\psi = c\psi$, so $I_{4}\psi = cJ^{-1}\psi = -cJ\psi$. Under the one metric, which by Proposition \claimref{dark-matter.thickness-not-gravitational-quantity} is not a field along the unshared directions, the rotation gauge takes values in shared bivectors, all of which commute with $J$; the content is then four-dimensional spin-$\tfrac12$ whose complex chirality (the action of $I_{4}$) is tied to the eigenvalue of the ungauged rotation $J$. In four dimensions the rotation-gauge anomaly form vanishes by (ii), and $J$ is not a gauge generator, so no obstruction arises from the one metric's field. If instead the tower's own gauge fields are fields on its six-dimensional space and act nontrivially on its constituents (the reading (G6), chosen knowing that the unit-ball fit needs residue $4$ empty; the store treats the five-dimensional towers' gauge fields this way, Proposition \claimref{dark-tower-matter.five-dimensional-towers}(i), (v), and the one rule's algebra acts on each constituent's pair factor, Theorem \claimref{synthesis.gauge-algebra-by-tower}), then by (ii) with $n = 6$ the gauge part $\tfrac{1}{24}\sum\mathrm{tr}_{R}F^{4}$ is positive on every acting generator, so the content's anomaly does not vanish and single-ideal matter of residue $4$ is inconsistent. Neither (G6) nor the cited propositions enter (i) to (iii); without (G6), (i) and (ii) leave residue $4$ unobstructed. The result filed alongside finds that with one ideal residue $4$ escapes emptiness under the momentum reading of mass ($m^{2} = \tfrac14$ for two constituents); under (G6) that escape is closed, since the single-ideal constituents are then anomalous whatever reading of mass is taken.

\emph{Consequence.} The six residues $1, 2, 3, 5, 6, 7$ meet no local anomaly from (ii): in odd dimension the forms have odd degree and vanish, and in dimension four the rotation-gauge part vanishes and the gauge part has odd power, so it can cancel within the content: a field and its conjugate of opposite chirality cancel identically, and the Standard Model content cancels by its hypercharges. So, in the lowest-causal window and given (A), under the single-ideal definition the anomaly conditions empty exactly residue $0$, using only the one metric and the sourcing reading, and residue $4$, using (G6). If both ideals are admitted at $n \equiv 2 \pmod 4$, the content contains each chirality with its conjugate in equal number, every form in (ii) cancels, and nothing here empties either residue; the content result then gives massive matter at $(1,1)$ as well.
